\documentclass[11pt]{article}
\usepackage[margin=0.9in]{geometry}
\usepackage{amsmath,amssymb,amsthm,mathtools}
\usepackage{graphicx}
\usepackage[colorlinks=true,linkcolor=blue,citecolor=blue,urlcolor=blue]{hyperref}
\usepackage{microtype}

\newtheorem{theorem}{Theorem}
\newtheorem{lemma}{Lemma}
\newtheorem{corollary}{Corollary}
\newcommand{\mal}{\operatorname{mal}}
\newcommand{\M}{\mathbb M}

\title{Malnormal contractions in every matrix dimension}
\author{Candidate full-proof packet}
\date{11 August 2026}

\begin{document}
\maketitle

\noindent\textbf{Status.}
Candidate full proof; likely valid; human review recommended.

\section{Source problem}

For $A\in\M_n$, Mulcahy and Sinclair define
\[
 \mal(A)=\inf\left\{\|[A,B]\|_2:
 B=B^*,\ \|B-\tau(B)I\|_2=1\right\},
\]
where $\tau$ is normalized trace and $\|\cdot\|_2$ is the unnormalized
Hilbert--Schmidt norm.  A sequence is malnormal when its operator norms are
uniformly bounded and its malnormality constants are uniformly positive.

The source constructs a universal $\kappa_0>0$ and contractions
$X_m\in\M_{3m}$ for every sufficiently large $m$, with
$\mal(X_m)\geq\kappa_0$~\cite{MS}.  It then asks for a sequence in every
dimension.  The exact statement (PDF page~12) is reproduced below.

\begin{center}
\includegraphics[width=0.94\textwidth]{figures/problem_crop.png}
\end{center}

We solve the problem by showing that malnormality can be padded upward one
dimension at a time with quantitative control.

\section{The padding lemma}

The key point is to add a new scalar block outside the spectrum and couple it
to the old space by one nonzero column.  Different blocks of a small
commutator then successively control the off-diagonal vector, the old
self-adjoint block modulo scalars, and the difference between the two scalar
parts.

\begin{lemma}[one-dimensional padding]\label{lem:padding}
Let $A\in\M_m$ satisfy $\|A\|\leq M$ and $\mal(A)\geq\kappa>0$.  Choose any
unit vector $v\in\mathbb C^m$, put $L=M+1$, and define
\[
 T=\begin{pmatrix}A&v\\0&L\end{pmatrix}\in\M_{m+1}.
\]
Then $\|T\|\leq M+2$ and
\[
 \mal(T)\geq D(\kappa,M)^{-1},
 \qquad
 D(\kappa,M)^2=\left(\frac2\kappa\right)^2+2+
 \left(2M+2+\frac2\kappa\right)^2.
\]
\end{lemma}

\begin{proof}
The norm estimate follows by splitting $T$ into its block diagonal and its
upper-right rank-one block.

Let
\[
 B=\begin{pmatrix}C&x\\x^*&b\end{pmatrix}=B^*,
 \qquad C\in\M_m,\quad x\in\mathbb C^m,\quad b\in\mathbb R,
\]
and set $\varepsilon=\|[T,B]\|_2$.  Direct multiplication gives
\[
 [T,B]=
 \begin{pmatrix}
 [A,C]+vx^* &(A-LI)x+vb-Cv\\
 x^*(LI-A)&-x^*v
 \end{pmatrix}.
\]
Each block has Hilbert--Schmidt norm at most $\varepsilon$.  Since
$L-\|A\|=1$, the lower-left block gives
\[
 \|x\|_2\leq\|x^*(LI-A)\|_2\leq\varepsilon.
\]
The upper-left block now yields
\[
 \|[A,C]\|_2\leq\varepsilon+\|vx^*\|_2\leq2\varepsilon.
\]
Writing $c=\tau(C)$ and using $\mal(A)\geq\kappa$, we obtain
\[
                    \|C-cI\|_2\leq\frac{2}{\kappa}\varepsilon.
\]
Finally, the upper-right block can be rewritten as
\[
 (A-LI)x+v(b-c)-(C-cI)v.
\]
Because $\|v\|=1$, it follows that
\begin{align*}
 |b-c|
 &\leq\varepsilon+(\|A\|+L)\|x\|_2+\|(C-cI)v\|_2\\
 &\leq\left(2M+2+\frac2\kappa\right)\varepsilon.
\end{align*}
Consequently,
\begin{align*}
 \|B-cI_{m+1}\|_2^2
 &=\|C-cI_m\|_2^2+2\|x\|_2^2+|b-c|^2\\
 &\leq D(\kappa,M)^2\varepsilon^2.
\end{align*}
The normalized trace scalar is the Hilbert--Schmidt best scalar
approximation, so
\[
 \|B-\tau(B)I_{m+1}\|_2\leq\|B-cI_{m+1}\|_2.
\]
This proves the malnormality estimate.
\end{proof}

\begin{corollary}[padding contractions]\label{cor:padding}
If $A$ is a contraction with $\mal(A)\geq\kappa$, then for any unit vector
$v$ the matrix
\[
             P_v(A)=\frac13\begin{pmatrix}A&v\\0&2\end{pmatrix}
\]
is a contraction in the next dimension and satisfies
\[
 \mal(P_v(A))\geq F(\kappa),\qquad
 F(\kappa)=\frac{1}{3D(\kappa,1)}>0.
\]
\end{corollary}

\begin{proof}
Apply Lemma~\ref{lem:padding} with $M=1$ and divide the resulting matrix by
three.  Both the operator norm and the commutator scale linearly.
\end{proof}

\section{A malnormal contraction in every dimension}

\begin{theorem}\label{thm:alln}
There are a universal constant $\kappa_*>0$ and contractions
$A_n\in\M_n$ for every $n\geq1$ such that
\[
       \|[A_n,B]\|_2\geq\kappa_*\|B-\tau(B)I_n\|_2
       \qquad(B=B^*\in\M_n).
\]
Thus Problem~5.1 of \cite{MS} has an affirmative answer.
\end{theorem}

\begin{proof}
Let $X_m\in\M_{3m}$ be the source's contractions for all $m\geq m_0$, with
$\mal(X_m)\geq\kappa_0$.  Put
\[
 \kappa_1=F(\kappa_0),\qquad
 \kappa_2=F(\kappa_1),\qquad
 \kappa_{\mathrm{large}}=\min\{\kappa_0,\kappa_1,\kappa_2\}>0.
\]
For $n=3m+r\geq3m_0$, where $r\in\{0,1,2\}$, start with $X_m$ and apply
Corollary~\ref{cor:padding} exactly $r$ times, choosing an arbitrary unit
padding vector at each step.  The resulting matrix lies in $\M_n$, is a
contraction, and has malnormality at least
$\kappa_{\mathrm{large}}$.

Only finitely many $n$ remain.  For $n\geq2$, let $S_n$ be the nilpotent
unilateral shift, a single Jordan block with ones on the superdiagonal.  Its
commutant consists of upper triangular Toeplitz matrices, and a self-adjoint
upper triangular matrix is diagonal.  Hence every self-adjoint matrix
commuting with $S_n$ is scalar.  Compactness of the unit sphere in the
finite-dimensional space of trace-zero self-adjoint matrices gives
$\mal(S_n)>0$.  For $n=1$ the required inequality is vacuous.  Taking the
minimum of $\kappa_{\mathrm{large}}$ and these finitely many positive constants
produces $\kappa_*>0$ and completes the sequence.
\end{proof}

\begin{corollary}[bounded-gap principle]
Suppose uniformly malnormal contractions exist in an unbounded set of matrix
dimensions whose successive gaps are bounded by $q$.  Then uniformly
malnormal contractions exist in every sufficiently large dimension.
\end{corollary}

\begin{proof}
Apply Corollary~\ref{cor:padding} at most $q-1$ times.  The finite list
$\kappa,F(\kappa),\ldots,F^{q-1}(\kappa)$ has a positive minimum.
\end{proof}

\section{Scope and literature boundary}

The construction inherits the nonconstructive choice of the source's quantum
expanders; it resolves the dimension gap, not the separate request for a
fully explicit algebraic seed.  It also does not address the sharp constant,
the random-matrix conjecture, or the proposed converse with quantum edge
expanders.

A bounded novelty search on 11 August 2026 used the exact title and problem
sentence, all-dimension malnormal sequences, one-dimensional block padding,
and subsequent references to the paper.  It found the source and unrelated
group-theoretic uses of ``malnormal,'' but no later solution of Problem~5.1 or
matching padding lemma.  Novelty confidence is moderate.

\paragraph{Human-review recommendation.}
Audit the four displayed blocks of $[T,B]$, especially the lower-left inverse
bound and the shared scalar $c=\tau(C)$.  Then check that the source's
Proposition~3.3 indeed supplies its $3m$ family for every sufficiently large
$m$, as used here.

\begin{thebibliography}{9}
\bibitem{MS}
G.~Mulcahy and T.~Sinclair,
\emph{Malnormal matrices},
Proc. Amer. Math. Soc. \textbf{150} (2022), 2989--3001;
arXiv:2009.11139.
\end{thebibliography}

\end{document}
